\section{Expanded reasoning: data partitions, loss derivatives and training choices}
\label{sec:expand-train}
\subsection{The 15\% holdout, seed and 20\% personal representation}
The inspected paths specify validation fraction 0.15 and split seed zero.
A seeded shuffle fixes which examples are withheld, provided the example
list and ordering are unchanged. The fraction controls data allocation;
the seed controls reproducibility, not statistical quality. For a constructed
100-line population a 15\% split with exact rounding gives 15 validation
and 85 training lines. Real per-writer rounding and minimum-set guards can
change totals. Never infer actual held-out counts from the fraction alone.

A line-level split can put neighbouring lines from one photographed page
in both sets. This tests new lines under partly shared capture conditions,
not new pages, pens or writers. A page-level split reduces that dependency;
a writer-disjoint split answers a different question from closed-set writer
identification. Fifteen percent is a sample-efficiency/precision compromise,
not a statistically universal optimum. After tuning on validation, it no
longer serves as a fully untouched final test.

The text trainer's personal fraction 0.20 controls representation during
mixed sampling~\cite{code2026}. It does not make public and personal data
equal: the intended shares are 80/20. In an ideal twenty-example batch,
that would mean sixteen public and four personal examples, subject to the
routine's rounding/replacement behaviour. Oversampling a small personal set
can repeat examples; an epoch need not then mean one visit per unique image.
This coefficient is separate from 0.2 dropout and from the 15\% holdout.
Its rationale is to prevent abundant public data overwhelming the personal
domain; it does not guarantee absence of bias or optimum transfer.

\subsection{Training settings: what each number controls}
The report's first-principles defaults use RMSprop for text, AdamW for
writers, initial learning rate $10^{-3}$, weight penalties $10^{-5}$ and
$10^{-4}$ respectively; text plateau reduction factor 0.5, patience 5 and
minimum $10^{-5}$; writer cosine annealing and gradient norm limit 5.
These describe code defaults, not a recovered configuration for every saved
checkpoint. Architecture sizes and preparation must also match the checkpoint.

RMSprop keeps $v_t=\alpha v_{t-1}+(1-\alpha)g_t^2$ and updates
$\theta\leftarrow\theta-\eta g_t/(\sqrt{v_t}+\epsilon)$.
This scales steps by recent squared-gradient magnitude~\cite{rmslecture}.
The inspected defaults are $\alpha=0.99,\epsilon=10^{-8}$. Larger
$\alpha$ gives a longer memory; the numerical floor prevents unstable small
denominators. The inspected RMSprop adds $\lambda\theta$ to the gradient,
so its weight penalty is coupled L2-style regularisation.

Adam forms $m_t=\beta_1m_{t-1}+(1-\beta_1)g_t$ and
$v_t=\beta_2v_{t-1}+(1-\beta_2)g_t^2$, corrects their initial bias using
$1-\beta_1^t$ and $1-\beta_2^t$, and scales the first moment by the square
root of the second~\cite{kingma2015}. Defaults 0.9 and 0.999 express
different smoothing times. AdamW additionally subtracts
$\eta\lambda\theta$ separately from the adaptive gradient update.
Decoupling matters because adding L2 to an adaptively scaled gradient is
not equivalent~\cite{loshchilov2019}. Neither optimiser proves a global
optimum for these nonconvex networks.

Halving learning rate after a plateau reduces update size when the monitored
validation quantity stops improving. ``Patience 5'' must be read with the
code's comparison: the inspected counter uses \texttt{badEpochs > patience},
so the reduction follows more than five counted bad epochs. Cosine decay
typically follows $\eta(t)=\eta_{\min}+
\tfrac12(\eta_0-\eta_{\min})[1+\cos(\pi t/T_{\max})]$ across the
configured run. These are schedules, not measured physical learning rates.
Changing loss reduction changes gradient scale and can invalidate a learning
rate comparison even if the displayed numerical rate is unchanged.

\subsection{Cross-entropy and its logit derivative from first principles}
For logits $z_k$, stable softmax is
$p_k=e^{z_k-m}/\sum_j e^{z_j-m}$ with $m=\max_jz_j$.
Subtracting the common $m$ leaves ratios unchanged. With true class $y$,
$L=-\log p_y=-z_y+\log\sum_j e^{z_j}$.
Direct differentiation gives
$\partial L/\partial z_k=p_k-\mathbf1[k=y]$.
For probabilities $(0.7,0.2,0.1)$ and true class one, loss is
$-\log0.7\simeq0.3567$ and gradient $(-0.3,0.2,0.1)$, summing to zero.
Gradient descent raises the true logit relative to the others. Class indices
are discrete labels; the derivative is with respect to the real logits.

A batch average $L_B=B^{-1}\sum_iL_i$ contributes a factor $1/B$ to every
per-example gradient. Apply that factor once, not at every backward layer.
An average makes nominal scale less sensitive to batch size, but the examples'
gradient directions still differ. Class imbalance and repeated samples can
still dominate the mean.

\subsection{CTC target occupancy, backward convention and batch scaling}
Use $Z=P(\ell\mid X)$ for total valid-path likelihood. Let
$\alpha_{t,s}$ include the emission at $(t,s)$ and let $\beta_{t,s}$
include only later emissions. Then conditional probability of state $s$ at
frame $t$ is $\exp(\alpha_{t,s}+\beta_{t,s}-\log Z)$ in log notation.
Summing over states carrying class $k$ gives
\[
q_{t,k}=\sum_{s:\ell'_s=k}
e^{\alpha_{t,s}+\beta_{t,s}-\log Z},\qquad
{\partial L\over\partial z_{t,k}}=p_t(k)-q_{t,k}.
\]
The condition $\ell'_s=k$ matters: state index and class index differ.
Different states can carry the same letter or blank. If backward scores
also include the current emission, divide it out once; otherwise it is
counted twice. This convention issue is a frequent source of incorrect
CTC implementations. The occupancy definition follows the valid-path
posterior, not the greedy path~\cite{gravesbook}.

For the constructed two-frame $a$ example above, at frame one valid-path
mass emitting $a$ is $0.42+0.18=0.60$ and blank mass is 0.28.
Thus $q_1(a)=0.60/0.88\simeq0.6818$, and the $a$ logit derivative is
$0.6-0.6818=-0.0818$. Occupancy sums to one at each valid frame; the logit
gradients sum to zero. These identities provide hand-checkable invariants.
An impossible target has $Z=0$ and requires explicit failure handling.

The project averages sequence losses across a batch without target-length
division. Equal coefficients $1/B$ do not mean equal influence per character:
long sequences can contribute larger loss/gradient. PyTorch's documented
\texttt{CTCLoss(reduction='mean')} divides by target length before averaging;
matching the project's batch mean requires summing per-sequence losses and
dividing by batch size~\cite{torchctc}. One must match blank index, valid
input/target lengths, tensor order and impossible-path behaviour before
calling a loss discrepancy a numerical bug.

\subsection{Backward pass: how every major layer receives a derivative}
For affine $Y=XW^T+b$ and upstream $G=\partial L/\partial Y$,
$dW=G^TX$, $dX=GW$, $db=\sum_{\rm batch}G$. A convolution can be
written as this matrix multiplication after stacking patches (im2col).
Its input gradient folds overlapping patch contributions back by \emph{adding}
them. Overwriting overlaps loses gradients. Bias sums over batch and spatial
positions. The padding added for forward computation must be removed from
the input gradient.

ReLU masks upstream values at nonpositive inputs. Max pooling scatters to
saved argmax locations, accumulating if needed. Average pooling distributes
equally. For batch normalisation, write $g_i=\partial L/\partial u_i$,
$r=(v+\epsilon)^{-1/2}$ and averages over the same $m$ values:
\[
d\beta=\sum_i g_i,\quad d\gamma=\sum_i g_i\hat z_i,\qquad
dz_i=\gamma r\,[g_i-\overline g-\hat z_i\overline{g\hat z}].
\]
This includes dependence of mean and variance on all batch values. Treating
them as constants gives the inference derivative, not the training derivative.
The distinction is mathematical; tolerances cannot repair a wrong formula.

For one LSTM step, combine upstream cell derivative $\bar c_t$ with the
hidden contribution $\bar h_t\odot o_t\odot(1-\tanh^2c_t)$.
Then $\bar f_t=\bar c_t\odot c_{t-1}$,
$\bar i_t=\bar c_t\odot g_t$,
$\bar g_t=\bar c_t\odot i_t$,
$\bar o_t=\bar h_t\odot\tanh c_t$ and
$\bar c_{t-1}=\bar c_t\odot f_t$ along the direct cell path.
Multiply gate derivatives by sigmoid/tanh derivatives, concatenate in the
same gate order as forward, form weight outer products, and propagate to
$x_t,h_{t-1}$. Reverse temporal order during backpropagation through time.
For the reverse-direction LSTM, undo the reversal when merging input
derivatives. For stacked layers, output derivatives split across directions
and then pass through the lower recurrent layer. These are explicit
chain-rule calculations, not an invocation of automatic differentiation.

\subsection{Clipping, initialisation, augmentation and their alternatives}
Global norm clipping uses
$g'=g\min(1,\tau/\|g\|_2)$ with $\tau=5$ for the inspected writer trainer.
A norm of 20 is scaled by 0.25; a norm of 2 is unchanged. It preserves
direction, unlike elementwise clipping. It limits a gradient norm, not
necessarily an adaptive optimiser's parameter-update norm. The method
addresses exploding gradients~\cite{pascanu2013}; threshold 5 remains a
project setting. The reported text route does not apply the same clipping.

Variance-scaled initialisation attempts to keep activations/gradients from
shrinking or growing through depth. For rectifiers, variance around
$2/n_{\rm in}$ motivates He-style scaling~\cite{he2015}; Xavier-style scaling
uses fan-in and fan-out for different activation assumptions~\cite{glorot2010}.
This does not make a freshly initialised model accurate. Learned batchnorm
scales starting at one and offsets at zero initially avoid an extra arbitrary
rescaling.

Augmentation perturbs only training examples with plausible shear, rotation,
perspective or elastic deformation~\cite{simard2003}. It assumes the label
is unchanged. An excessive warp can violate that assumption or create
nonphysical strokes. Validation preparation remains fixed, with dropout
disabled and saved normalisation statistics, so its score is not an average
over random training perturbations. Stronger augmentation is not automatically
better, especially for writer identity where style cues themselves are labels.

\subsection{Extended training algorithm and historical time evidence}
Conceptually: choose fixed partitions; sample a batch and any domain mixture;
augment only training images; apply the checkpoint's preparation; run the
train-mode backbone and task head while saving forward caches; compute CTC
or cross-entropy; apply the correct batch reduction; backpropagate through
the head, pooling and convolution stack; clip if configured; update optimiser
state/weights; at epoch end validate deterministically; update schedule;
save a progress checkpoint and replace best only under the selected criterion.
Text selection uses validation loss; writer selection uses validation accuracy.
Resume requires optimiser moments, epoch/scheduler state and preferably
random state, not just weights, to reproduce a continuous run.

The final report's Appendix C.6 retains writer epoch times around 19.5--20.9
minutes and selected reader epochs 7, 8, 11 and 16 at 122.3, 120.1, 111.6
and 112.3 minutes. These are reported log rows, not timings measured here.
Multiplying one selected epoch by a desired epoch count estimates neither
the exact complete run time nor the time on another host. Caches, validation,
batch sizes, implementation and interruptions matter. ``Best saved epoch''
is not necessarily final epoch, and one loss value is not an accuracy.
The logs and checkpoint metadata are needed to tie any quoted result to
its host, precision, training route and selected model. No new training,
inference or numerical comparison was run for this edition.
